% Generated draft: requires mathematical and source review.
\documentclass[11pt]{article}
\usepackage[margin=25mm]{geometry}
\usepackage{fontspec}
\setmainfont{Latin Modern Roman}
\usepackage{amsmath,amssymb,mathtools,unicode-math}
\setmathfont{Latin Modern Math}
\usepackage{enumitem,needspace,xurl,hyperref}
\setlength{\parindent}{0pt}
\setlength{\parskip}{5pt}
\setlength{\emergencystretch}{3em}
\newcommand{\N}{\mathbb{N}}
\newcommand{\Z}{\mathbb{Z}}
\newcommand{\Q}{\mathbb{Q}}
\newcommand{\R}{\mathbb{R}}
\newcommand{\C}{\mathbb{C}}
\newcommand{\F}{\mathbb{F}}
\newcommand{\A}{\mathbb{A}}
\newcommand{\PP}{\mathbb{P}}
\newcommand{\E}{\mathbb{E}}
\newcommand{\eps}{\varepsilon}
\newcommand{\dd}{\,\mathrm{d}}
\begin{document}
\section*{UnsolvedMath ID: 700805}
\textbf{Source problem code:} catalogue entry
\subsection{Definitions and mathematical statement}
A group \(G\) is periodic if every element of \(G\) has finite order. Define its lower central series by \(\gamma_1(G)=G\) and \(\gamma_{i+1}(G)=[\gamma_i(G),G]\); the group is nilpotent if \(\gamma_{c+1}(G)=1\) for some finite integer \(c\geq 0\). Here \([H,K]\) denotes the subgroup generated by the commutators \(h^{-1}k^{-1}hk\), with \(h\in H\) and \(k\in K\). Define the derived series by \(G^{(0)}=G\) and \(G^{(r+1)}=[G^{(r)},G^{(r)}]\). Its derived length is \(\operatorname{dl}(G)=\min\{r\geq 0:G^{(r)}=1\}\).

Interpret a regular automorphism \(\alpha\in\operatorname{Aut}(G)\) as a fixed-point-free automorphism, meaning
\[
C_G(\alpha)=\{g\in G:\alpha(g)=g\}=\{1\}.
\]
The supplied record does not explicitly define ``regular'', so this interpretation requires source confirmation.

Does there exist a function
\[
f:\{p\in\Z_{\geq 2}:p\text{ is prime}\}\times\Z_{\geq 1}\longrightarrow\Z_{\geq 0}
\]
such that, for every prime \(p\), every integer \(n\geq 1\), every nilpotent periodic group \(G\), and every regular automorphism \(\alpha\) of \(G\) of exact order \(p^n\), one has
\[
\operatorname{dl}(G)\leq f(p,n)?
\]
The bound is to be independent of \(G\), including its nilpotency class. The source imposes no finiteness assumption on \(G\) and no additional coprimality hypothesis.
\subsection{Short English statement}
Must the derived length of every nilpotent periodic group admitting a regular automorphism of order \(p^n\) be bounded solely in terms of the prime \(p\) and the positive integer \(n\)? Here regular is interpreted as fixed-point-free, subject to confirmation of the source terminology.
\subsection{Sources}
\begin{itemize}
\item Local UnsolvedMath registry, permanent ID \(700805\), exact source code \texttt{catalogue entry}, title ``Kourovka Notebook Problem 11.113''. The supplied record labels the problem open; this is a supplied status, not an independently verified mathematical or literature assessment.
\item \emph{The Kourovka Notebook (Unsolved Problems in Group Theory)}, original problem 11.113, 11th issue (1990); the supplied background also identifies the 21st issue edition (2026), arXiv:1401.0300v48. The collection description names E. I. Khukhro and V. D. Mazurov as editors. This reference is supplied by the record and has not been independently consulted.
\item Attribution is unresolved within the record: the statement begins ``(B. Hartley)'', whereas the background and proposer field identify E. I. Khukhro, with proposed year 1990. Both attributions are retained without assigning unprovided roles.
\end{itemize}
\end{document}
